\documentclass[11pt]{article}
\usepackage{fontspec,amsmath,amssymb,amsthm,hyperref}
\usepackage[margin=1in]{geometry}
\setlength{\emergencystretch}{3em}
\title{Compactly supported smooth tests determine a finite measure}
\author{LangevinLean mathematical exposition}
\date{}
\begin{document}\maketitle
\noindent Source module: \texttt{UnadjustedLangevin/TestMeasure.lean}. Compile this proof document with XeLaTeX or LuaLaTeX. The accompanying HTML places the complete Lean declarations beside every proof.


On a finite-dimensional real normed space, finite measures agreeing on all compactly supported smooth tests are equal. Agreement on compactly supported \(C^2\) tests is sufficient as well.


\section{Compare Radon–Nikodym densities}


Let \(\xi=\mu+\nu\), so both measures are absolutely continuous with respect to \(\xi\). Their finite Radon–Nikodym densities are locally integrable. Rewrite the given test integrals against these densities; their difference integrates to zero against every compactly supported smooth test. The fundamental lemma for test functions gives equality of the real-valued densities \(\xi\)-almost everywhere. Finiteness allows conversion back to extended densities, and the Radon–Nikodym reconstruction formulas give \(\mu=\nu\). The \(C^2\) version applies this result because every smooth test is \(C^2\).



\noindent Lean declarations: \texttt{measure\_eq\_of\_integral\_contDiff\_eq}, \texttt{measure\_eq\_of\_integral\_contDiff\_two\_eq}.

\end{document}
